\section{The Problem We Are Actually Solving}
\label{sec:problem}

\paragraph{The customer's question.}
A chip team with a fully closed-source design wants an internal assistant that
knows the chip the way a senior engineer does. Before committing to a training
program, two prior questions must be answered with evidence rather than
enthusiasm.

\paragraph{Question 1: where can a self-trained model beat the frontier?}
Not on public skills. We reproduced the VeriThoughts fine-tuning
pipeline~\citep{yubeaton2025verithoughts} as a calibration exercise and then
ran Claude Opus~5 on the same formally-verified Verilog benchmark: the frontier
model scores 95.3\% pass@1, above the fine-tuned specialist models the
literature celebrates. Public RTL generation is saturating; a 7--9B model
fine-tuned on public data differentiates nothing. What a frontier model
\emph{cannot} have is knowledge that never reached its training set---the
proprietary chip itself. Knowledge injection is therefore not one use case
among many; it is the only battleground where the self-trained model holds
structural advantage.

\paragraph{Question 2: what counts as ``knowing''?}
We require \emph{closed-book} recall: no retrieval, no context documents, the
fact must come out of the weights. This is a deliberate, restrictive choice.
Retrieval-augmented generation answers lookup questions well, but the
engineer-facing failures we target---cross-file reasoning, ``which function
configures this register,'' recall inside a debugging conversation where
nobody pastes the register map---degrade precisely when knowledge lives only
in an external index. Testing closed-book recall isolates the property we are
trying to buy with training compute. (Whether RAG plus a frontier model is the
better \emph{product} for a given team is a separate question our open
questions section returns to; this study measures what training can inject.)

\paragraph{Requirements.}
The study therefore needs: (i) a testbed whose facts are verifiably absent or
scarce in frontier pretraining, yet real enough that results transfer to a
customer engagement; (ii) an evaluation that is machine-checkable, layered
across the kinds of knowledge an engineer uses, and hardened against the
model answering by surface pattern rather than knowledge; (iii) a recipe
whose every component is justified by an ablation, because the customer will
ask ``do we need this step?''\ for each one; and (iv) a footprint a small
team can afford without procurement: every training run in this study fits on
a single 4-GPU node in about an hour.
